\chapter{Fundamentação Teórica}
\label{cap:fundamentacao}

Este capítulo apresenta a base conceitual da plataforma. A primeira parte trata do lado técnico: a arquitetura de sistemas operacionais, o padrão POSIX e a organização do sistema de arquivos no Linux. A segunda parte trata do lado pedagógico: a Teoria da Carga Cognitiva, a Aprendizagem Ativa, o \textit{Microlearning} e a dispersão causada pelos \textit{smartphones} em sala de aula.

% ----------------------------------------------------------
\section{Arquitetura de Sistemas Operacionais e o Padrão POSIX}
\label{sec:posix}
% ----------------------------------------------------------

Um sistema operacional é a camada que fica entre o \textit{hardware} e os programas do usuário \cite{silberschatz2018operating, tanenbaum2015modern}. Segundo \citeonline{tanenbaum2015modern}, ele cumpre dois papéis. O primeiro é o de \textbf{gerenciador de recursos}: distribui processador, memória, discos e dispositivos de entrada e saída entre programas que competem por eles. O segundo é o de \textbf{máquina estendida} (\textit{extended machine}): oferece abstrações convenientes, como arquivos, processos e \textit{sockets}, que escondem os detalhes do \textit{hardware}.

Os programas em espaço de usuário (\textit{user space}) conversam com o núcleo (\textit{kernel space}) por meio das chamadas de sistema (\textit{system calls} ou \textit{syscalls}). Com a multiplicação de variantes proprietárias do Unix, essas interfaces divergiram e a compatibilidade entre sistemas se perdeu. Para resolver o problema, o IEEE, a ISO e \textit{The Open Group} criaram o padrão POSIX (\textit{Portable Operating System Interface}) \cite{posix2018}.

O POSIX define contratos uniformes para operações de Entrada e Saída baseadas em descritores de arquivo, para criação e controle de processos (\texttt{fork}, \texttt{exec}, \texttt{wait}), para sinais assíncronos e para o comportamento do interpretador de comandos (\textit{shell}) e dos utilitários de linha de comando \cite{kerrisk2010linux}. Quem domina essas interfaces consegue escrever programas portáveis entre as diferentes distribuições Linux e os demais sistemas Unix-like.

% ----------------------------------------------------------
\section{Sistemas de Arquivos e o Padrão FHS}
\label{sec:fhs}
% ----------------------------------------------------------

No Unix e no Linux vale o princípio de que ``tudo é um arquivo'' (\textit{everything is a file}). Não só os dados em disco, mas também dispositivos, periféricos, canais de comunicação entre processos e estruturas do próprio núcleo aparecem como nós de uma única árvore de diretórios, com raiz em \texttt{/} \cite{nemeth2017unix, kerrisk2010linux}.

Para que cada distribuição não organize essa árvore à sua maneira, a \textit{Linux Foundation} mantém o \textit{Filesystem Hierarchy Standard} (FHS) \cite{fhs30}. O Quadro~\ref{qua:fhs} resume os principais diretórios definidos pelo padrão.

\begin{quadro}[htbp]
    \caption{Principais diretórios definidos pelo FHS}
    \label{qua:fhs}
    \footnotesize
    \begin{tabularx}{\textwidth}{|p{3.4cm}|L|p{3.6cm}|}
        \hline
        \textbf{Diretório} & \textbf{Finalidade} & \textbf{Exemplo} \\
        \hline
        \texttt{/bin}, \texttt{/sbin} & Binários essenciais para operar e administrar o sistema & \texttt{ls}, \texttt{mount} \\
        \hline
        \texttt{/etc} & Configurações globais da máquina e dos serviços & \texttt{/etc/passwd} \\
        \hline
        \texttt{/home} & Dados pessoais dos usuários comuns & \texttt{/home/ana} \\
        \hline
        \texttt{/root} & Diretório pessoal do superusuário & --- \\
        \hline
        \texttt{/var} & Dados que mudam durante a operação: \textit{logs}, bancos de dados, filas & \texttt{/var/log/syslog} \\
        \hline
        \texttt{/opt}, \texttt{/usr/local} & \textit{Software} adicional e programas compilados localmente & \texttt{/usr/local/bin} \\
        \hline
        \texttt{/tmp} & Arquivos temporários, protegidos pelo \textit{sticky bit} & \texttt{/tmp/sessao.lock} \\
        \hline
        \texttt{/dev}, \texttt{/proc}, \texttt{/sys} & Sistemas de arquivos virtuais que expõem dispositivos e estruturas do \textit{kernel} & \texttt{/dev/sda}, \texttt{/proc/cpuinfo} \\
        \hline
    \end{tabularx}
    \fonte{Adaptado de \citeonline{fhs30}.}
\end{quadro}

A segurança do sistema de arquivos se apoia em identidades de usuário (\textit{User ID} -- UID), em grupos primários e secundários (\textit{Group ID} -- GID) e na matriz de permissões POSIX: leitura (\textit{r} = 4), escrita (\textit{w} = 2) e execução (\textit{x} = 1), aplicadas ao dono (\textit{user}), ao grupo (\textit{group}) e aos demais (\textit{others}) \cite{nemeth2017unix}. Há ainda atributos especiais como \textit{SetUID}, \textit{SetGID} e o \textit{sticky bit} (octal 1000). Este último é o que impede, em diretórios compartilhados como \texttt{/tmp}, que um usuário apague arquivos criados por outro.

% ----------------------------------------------------------
\section{A Teoria da Carga Cognitiva no Ensino de Computação}
\label{sec:carga-cognitiva}
% ----------------------------------------------------------

A Teoria da Carga Cognitiva (\textit{Cognitive Load Theory} -- CLT), proposta por \citeonline{sweller1988cognitive} e ampliada nas décadas seguintes \cite{sweller2011cognitive, paas2003cognitive}, parte da distinção entre memória de trabalho (\textit{working memory}) e memória de longo prazo (\textit{long-term memory}).

A memória de longo prazo tem capacidade praticamente ilimitada para guardar estruturas de conhecimento, os chamados esquemas mentais. A memória de trabalho, ao contrário, é um gargalo: consegue manter e manipular poucos elementos novos ao mesmo tempo, e por pouco tempo \cite{sweller1988cognitive}. Quando uma atividade exige mais do que ela suporta, ocorre sobrecarga cognitiva (\textit{cognitive overload}) e o aprendizado não acontece.

A literatura da área \cite{paas2003cognitive, sweller2011cognitive} divide o esforço mental em três tipos de carga, descritos no Quadro~\ref{qua:cargas-cognitivas} com exemplos do ensino de Sistemas Operacionais.

\begin{quadro}[htbp]
    \caption{Os três tipos de carga cognitiva aplicados ao ensino de SO}
    \label{qua:cargas-cognitivas}
    \footnotesize
    \begin{tabularx}{\textwidth}{|p{2.4cm}|L|L|}
        \hline
        \textbf{Carga} & \textbf{O que é} & \textbf{Exemplo no ensino de SO} \\
        \hline
        \textbf{Intrínseca} & Complexidade do próprio conteúdo e do número de elementos que interagem entre si & Entender a precedência das permissões octais, a sincronização de processos ou o encadeamento de \textit{pipes} \\
        \hline
        \textbf{Germana} & Esforço produtivo de construir e consolidar esquemas na memória de longo prazo & Relacionar \texttt{chmod}, \texttt{umask} e \texttt{ls -l} em um modelo mental único de permissões \\
        \hline
        \textbf{Extrínseca} & Esforço desperdiçado com a forma de apresentação ou com obstáculos que não ensinam nada & Configurar o hipervisor, ativar VT-x na BIOS, depurar a rede da máquina virtual \\
        \hline
    \end{tabularx}
    \fonte{Autoria própria (2026), com base em \citeonline{paas2003cognitive} e \citeonline{sweller2011cognitive}.}
\end{quadro}

No ensino prático tradicional, pedir a iniciantes que instalem hipervisores, ajustem a BIOS/UEFI e consertem a rede da VM gera uma grande \textbf{carga extrínseca} \cite{silva2021ensino}. Essa carga ocupa a memória de trabalho e esgota a atenção do aluno antes que ele chegue aos conceitos de sistemas operacionais propriamente ditos.

Por isso, uma ferramenta educacional deve começar eliminando a carga extrínseca, oferecendo um ambiente que esteja pronto para uso imediato \cite{guo2013online}.

% ----------------------------------------------------------
\section{Aprendizagem Ativa, Microlearning e Ambientes Web Interativos}
\label{sec:aprendizagem-ativa}
% ----------------------------------------------------------

O ensino expositivo, baseado na transmissão passiva de conteúdo e em longas apostilas em PDF, tem eficácia limitada para desenvolver habilidades práticas em computação \cite{benari2001constructivism, freeman2014active}.

Em uma meta-análise publicada nos \textit{Proceedings of the National Academy of Sciences}, \citeonline{freeman2014active} compararam centenas de estudos e mostraram que a Aprendizagem Ativa (\textit{Active Learning}) em disciplinas de ciências, tecnologia, engenharia e matemática (STEM) eleva as notas médias e reduz em 55\% as taxas de reprovação em relação às aulas puramente expositivas.

O \textit{Microlearning}, por sua vez, propõe dividir conteúdos complexos em pequenas unidades (\textit{micro-lessons}), cada uma com um único objetivo prático e imediato \cite{kapp2019microlearning}. Cada lição combina uma explicação curta com um desafio que se resolve na hora, criando ciclos rápidos de tentativa, erro e acerto.

Ambientes interativos que rodam no próprio navegador, sem instalar nada, têm se destacado na literatura de Educação em Computação (SIGCSE/ACM) como forma de aumentar o engajamento e democratizar o acesso \cite{guo2013online, silva2021ensino}. Como mostra \citeonline{guo2013online}, remover a etapa de instalação de compiladores e interpretadores reduz drasticamente o atrito inicial e aumenta o tempo dedicado à experimentação.

% ----------------------------------------------------------
\section{O Perfil dos Novos Discentes e a Dispersão Atencional por Smartphones}
\label{sec:dispersao-smartphones}
% ----------------------------------------------------------

Quem ingressa no ensino superior nesta década pertence a uma geração cujos hábitos de atenção foram moldados pelos \textit{smartphones} e pela conexão permanente à internet \cite{kuznekoff2013impact}. Esse público consome informação de forma fragmentada, espera respostas imediatas dos sistemas e tende à multitarefa digital.

Estudos sobre o uso de celulares em sala de aula no ensino superior mostram que a distração digital prejudica a retenção de conteúdo e o engajamento, sobretudo quando a atividade perde dinamismo \cite{kuznekoff2013impact}. Em um laboratório, o atrito técnico excessivo, como esperar a máquina virtual iniciar ou encarar uma tela congelada, funciona como gatilho: o estudante deixa o computador de lado e pega o celular.

Depois que a atenção migra para redes sociais e mensagens, recuperá-la custa caro, e o aluno perde o ritmo da aula \cite{kuznekoff2013impact, silva2021ensino}. Ambientes com resposta visual instantânea e terminal interativo no navegador fazem o caminho inverso: prendem a atenção do estudante na atividade desde os primeiros minutos.
